% Part: incompleteness
% Chapter: introduction
% Section: overview

\documentclass[../../../include/open-logic-section]{subfiles}

\begin{document}

\olfileid{inc}{int}{ovr}

\olsection{نظرة عامة على نتائج عدم الاكتمال}

كان هيلبرت يرجو أن تجمع الرياضيات في نظرية
!!{axiomatizable} يثبت أنها !!{complete} و!!{decidable}،
وأن يثبت اتساقها بوسائل «نهائية» ضعيفة جدًا؛ فيكون
ذلك جوابًا عن اعتراضات الحدسية على الرياضيات
الكلاسيكية. وجاءت مبرهنتا غودل لتبيّنا تعذر بلوغ
هذه المطالب مجتمعة.

أثبتت مبرهنة غودل الأولى أن إحدى صور
\emph{مبادئ الرياضيات} لراسل ووايتهد ليست
!!{complete}. ولم يكن البرهان خاصًا بها، بل
تناول في حقيقته طائفة واسعة من النظريات.
فلم يكن قصور \emph{مبادئ الرياضيات} عن استيعاب
الرياضيات تمامًا عيبًا يزول بمجرد استبدالها؛
إذ \emph{لا} نظرية مستوفية للشروط المطلوبة
تبلغ ذلك. وقد احتاج تعيين تلك الشروط وتحرير
البرهان حتى لا يعتمد على غيرها إلى وقت.
أما الآن فلنا صيغة عامة جدًا للمبرهنة الأولى
في لغة الحساب $\Lang L_{\OLId{latin-looped}{A}}$.

\begin{thm}
لتكن ${\OLId{greek-variable}{Gamma}}$ نظرية متسقة !!{axiomatizable}
في~$\Lang L_{\OLId{latin-looped}{A}}$، ولنفترض أنها !!{represents} جميع الدوال
القابلة للحساب والعلاقات !!{decidable}.
عندئذ لا تكون ${\OLId{greek-variable}{Gamma}}$ !!{complete}.
\end{thm}
% تصحيح صياغة كلاسيكية (C278-01): أُسند الفعل إلى Gamma صراحةً بدل التركيب المتنافر «ولتكن تمثّل».

ومعنى أن ${\OLId{greek-variable}{Gamma}}$ ليست !!{complete} أن ثمة
!!{sentence}~$!A$ واحدة على الأقل تحقق
${\OLId{greek-variable}{Gamma}}\Proves/ !A$ و${\OLId{greek-variable}{Gamma}}\Proves/\lnot !A$.
وتسمى !!a{sentence} بهذه الصفة \emph{مستقلة}
عن~${\OLId{greek-variable}{Gamma}}$. وإثبات وجود جمل مستقلة ممكن
بحجة قصيرة نسبيًا؛ لكن مزية برهان غودل
تعيين !!{sentence} مستقلة بعينها. أما
المبرهنة بصيغتها السابقة، تحت الاتساق وحده،
فمأخذها تحسين روسر لحجة غودل. وفي البناء
الأصلي نحصل على !!a{sentence} معينة
$!G_{\OLId{greek-variable}{Gamma}}$، تسمى \emph{جملة غودل}
لـ~${\OLId{greek-variable}{Gamma}}$، وتكافئ في~${\OLId{greek-variable}{Gamma}}$ القول
بأن $!G_{\OLId{greek-variable}{Gamma}}$ غير قابلة للبرهنة في~${\OLId{greek-variable}{Gamma}}$.
فإذا اتسقت ${\OLId{greek-variable}{Gamma}}$ امتنعت عليها برهنة
$!G_{\OLId{greek-variable}{Gamma}}$؛ ولإثبات استقلال $!G_{\OLId{greek-variable}{Gamma}}$
بالحجة الأصلية، أي عجز ${\OLId{greek-variable}{Gamma}}$ عن برهنة
نفيها أيضًا، نحتاج إلى فرض أقوى، كاتساق~$\omega$.
والبناء \emph{إنشائي}: فمن أكسمة~${\OLId{greek-variable}{Gamma}}$
ووصف نظام ال!!{derivation} نستخرج طريقة
نكتب بها~$!G_{\OLId{greek-variable}{Gamma}}$ فعلًا.

ويقتضي هذا البناء أن نعبر في~$\Lang L_{\OLId{latin-looped}{A}}$
عن خواص الحدود وال!!{formula}s والعمليات
عليها في~$\Lang L_{\OLId{latin-looped}{A}}$ نفسها. فمن الخواص:
«$!A$ !!a{sentence}»، و«$\delta$
!!a{derivation} لـ~$!A$»؛ ومن العمليات
$\Subst{!A}{{\OLId{latin-ordinary}{t}}}{{\OLId{latin-ordinary}{x}}}$. ولهذا التمثيل مطلبان.
أولهما أن ترمز الرموز ومتتالياتها، أي الحدود
وال!!{formula}s وال!!{derivation}s ونحوها،
بأعداد طبيعية، فهي موضوع~$\Lang L_{\OLId{latin-looped}{A}}$.
وثانيهما أن يكون الترميز بحيث تثبت ${\OLId{greek-variable}{Gamma}}$
الحقائق المطلوبة. فنبيّن أن الخواص تصير
خواص !!{decidable} للأعداد، وأن العمليات
تصير دوال قابلة للحساب عليها. وهذا ما يسمى
«حوسبة التركيب».

وقبل تفصيل حوسبة التركيب نبحث شرط القوة
في~${\OLId{greek-variable}{Gamma}}$: أن تكون !!{represents} جميع
الدوال القابلة للحساب والعلاقات !!{decidable}.
ولا بد لذلك من حد دقيق لقابلية الحساب.
وله حدود كثيرة، كالحساب بآلات تورنغ أو
ببرامج لغة عامة الغرض. لكن غرضنا أن
!!{represents} نظرية في~$\Lang L_{\OLId{latin-looped}{A}}$
تلك الدوال والعلاقات؛ فنختار حدًا بسيطًا
يخص الدوال العددية، هو مفهوم \emph{الدالة
العودية}. فنبحث هذه الدوال أولًا، ثم نثبت
أن $\Th{Q}$ !!{represents} الدوال والعلاقات
العودية جميعًا. وبذلك يصح تطبيق مبرهنة
عدم الاكتمال على $\Th{Q}$ و$\Th{PA}$،
بعد أن نثبت كفايتهما من جهة القوة.

وتنتهي حوسبة التركيب إلى !!a{formula}
$\OProv[{\OLId{greek-variable}{Gamma}}]({\OLId{latin-ordinary}{x}})$ تعبّر عن قابلية البرهنة
بترميز ال!!{formula}s أعدادًا. وهي تتناول
البرهنة من بديهيات~${\OLId{greek-variable}{Gamma}}$: فإذا كان
رمز $!A$ العدد~${\OLId{latin-ordinary}{n}}$ وكان ${\OLId{greek-variable}{Gamma}}\Proves !A$،
ثبت ${\OLId{greek-variable}{Gamma}}\Proves\Prov[{\OLId{greek-variable}{Gamma}}](\num{{\OLId{latin-ordinary}{n}}})$.
وبهذا «المحمول لقابلية البرهنة» الخاص
بـ${\OLId{greek-variable}{Gamma}}$، نستطيع أيضًا التعبير عن اتساق
${\OLId{greek-variable}{Gamma}}$، بمعنى معين، ب!!a{sentence}
من~$\Lang L_{\OLId{latin-looped}{A}}$. فعبارة الاتساق
لـ~${\OLId{greek-variable}{Gamma}}$ هي ال!!{sentence}
$\lnot\Prov[{\OLId{greek-variable}{Gamma}}](\num{{\OLId{latin-ordinary}{n}}})$، حيث
${\OLId{latin-ordinary}{n}}$ رمز تناقض، كـ~$\lfalse$.
ومفاد المبرهنة الثانية أن النظريات المتسقة
!!{axiomatizable}، متى استوفت شروطها،
لا تثبت عبارات اتساقها. وهذه الشروط
أشد قليلًا من تمثيل الدوال القابلة
للحساب والعلاقات !!{decidable} جميعًا؛
وسنثبت تحققها في $\Th{PA}$.

\end{document}
